% ============================================================
% Sections 4.2 and 4.3: Full-variance optimal weights and
% the Two-Step Estimator
%
% These sections follow Section 4.1 (A1-optimal weights,
% Lagrange multiplier argument) and precede Section 5
% (simulation study).
%
% Notation follows Yang Wu Azzollini DPhil thesis, Chapter 5.
% ============================================================

% ----------------------------------------------------------
\subsection{Full-variance optimal weights}\label{sec:fullopt}
% ----------------------------------------------------------

The weights $w_{ij} = \tau_{ij}/(\tau_i^X\tau_j^Y)$ minimise the
coefficient $A_1$ of $\sigma_1^2\sigma_2^2$ in Theorem~\ref{thm:genHY}.
A natural question is whether these weights also minimise the full
conditional variance
\begin{equation}\label{eq:fullvar}
    V(\mathbf{w}) \;=\; \sigma_1^2\sigma_2^2\, A_1(\mathbf{w})
                      \;+\; \sigma_{12}^2\, A_2(\mathbf{w}).
\end{equation}
The answer is \emph{no} in general: as we now show, the
full-variance optimum depends on the unknown ratio
$\kappa = \sigma_{12}^2/(\sigma_1^2\sigma_2^2) = \rho^2$,
and it coincides with the $A_1$-optimal weights only when $\rho = 0$.

\paragraph{Homogeneity.}
Both $A_1(\mathbf{w})$ and $A_2(\mathbf{w})$ are homogeneous of
degree zero in $\mathbf{w}$: multiplying every $w_{ij}$ by a positive
constant leaves $A_1$ and $A_2$ unchanged, since numerator and
denominator in \eqref{eq:genZHY} are both quadratic in $\mathbf{w}$.
We may therefore impose the normalisation constraint
\begin{equation}\label{eq:normalise}
    T(\mathbf{w}) \;\equiv\; \sum_{ij} \tau_{ij}\, w_{ij} \;=\; 1
\end{equation}
without loss of generality.

\paragraph{The optimisation problem.}
Under constraint~\eqref{eq:normalise} the objective
$V(\mathbf{w})$ is a \emph{strictly convex quadratic} in
$\mathbf{w}$, since both $A_1$ and $A_2$ become convex quadratics
once the denominator $T(\mathbf{w})^2$ is fixed at~$1$.
It therefore has a unique global minimum, characterised
by the Karush--Kuhn--Tucker (KKT) conditions.

\begin{thm}[Full-variance optimal weights]\label{thm:kkt}
Under the normalisation~\eqref{eq:normalise}, the weights
$\mathbf{w}^*(\kappa)$ minimising the full conditional variance
$V(\mathbf{w}) = \sigma_1^2\sigma_2^2 A_1(\mathbf{w})
              + \sigma_{12}^2 A_2(\mathbf{w})$
are the unique solution of the linear system
\begin{equation}\label{eq:kkt}
    w_{ij}^*\bigl[\tau_i^X \tau_j^Y - \kappa\,\tau_{ij}^2\bigr]
    + \kappa\,\tau_{ij}\bigl[R_i(\mathbf{w}^*) + C_j(\mathbf{w}^*)\bigr]
    = \mu\,\tau_{ij},
    \qquad \tau_{ij} > 0,
\end{equation}
together with constraint~\eqref{eq:normalise}, where
\begin{align*}
    R_i(\mathbf{w}) &= \sum_j w_{ij}\,\tau_{ij},\quad
    C_j(\mathbf{w}) = \sum_i w_{ij}\,\tau_{ij},
\end{align*}
are the row and column sums of the matrix $\{w_{ij}\tau_{ij}\}$,
$\kappa = \sigma_{12}^2/(\sigma_1^2\sigma_2^2)$,
and $\mu$ is a Lagrange multiplier for
constraint~\eqref{eq:normalise}.

The system~\eqref{eq:kkt} is linear in the unknowns
$\{w_{ij}^*: \tau_{ij}>0\}$ and $\mu$, with dimension
$(K+1)\times(K+1)$ where $K = |\{(i,j): \tau_{ij}>0\}|$
is the number of overlapping return-period pairs.
\end{thm}

\begin{proof}
Under normalisation~\eqref{eq:normalise} the objective becomes
\begin{align*}
    V(\mathbf{w}) &= \sigma_1^2\sigma_2^2
    \sum_{ij} w_{ij}^2\,\tau_i^X\tau_j^Y
    + \sigma_{12}^2
    \Bigl[\sum_i R_i(\mathbf{w})^2
         + \sum_j C_j(\mathbf{w})^2
         - \sum_{ij} w_{ij}^2\,\tau_{ij}^2\Bigr].
\end{align*}
Differentiating with respect to $w_{ij}$ (for $\tau_{ij}>0$)
and introducing Lagrange multiplier $\tilde\mu$ for
constraint~\eqref{eq:normalise} gives the stationarity condition
\begin{align*}
    2\sigma_1^2\sigma_2^2\,w_{ij}\,\tau_i^X\tau_j^Y
    + 2\sigma_{12}^2\,\tau_{ij}\bigl[R_i(\mathbf{w}) + C_j(\mathbf{w})
       - w_{ij}\,\tau_{ij}\bigr]
    &= \tilde\mu\,\tau_{ij}.
\end{align*}
Dividing through by $2\sigma_1^2\sigma_2^2$ and setting
$\mu = \tilde\mu/(2\sigma_1^2\sigma_2^2)$ yields~\eqref{eq:kkt}.
Since the objective is strictly convex and the constraint is
affine, the KKT system has a unique solution.
\end{proof}

\begin{remark}[Reduction to $A_1$-optimal weights at $\rho=0$]
When $\kappa = 0$, i.e.\ $\rho = 0$, the KKT
system~\eqref{eq:kkt} reduces to
$w_{ij}^*\,\tau_i^X\tau_j^Y = \mu\,\tau_{ij}$,
giving $w_{ij}^* \propto \tau_{ij}/(\tau_i^X\tau_j^Y)$, which is
exactly the $A_1$-optimal weight. This confirms that the $A_1$-optimal
weights are also full-variance optimal when $\rho=0$.
\end{remark}

\begin{remark}[Dependence on the unknown $\kappa$]
The system~\eqref{eq:kkt} depends on $\kappa = \rho^2$, which is
itself a function of the parameter $\sigma_{12}$ being estimated.
The full-variance optimal weights are therefore \emph{infeasible}:
they cannot be computed directly from the data without knowledge of
$\rho$. This is the direct analogue of the generalised least squares
(GLS) problem, where the optimal weighting matrix depends on the
unknown covariance structure.

The magnitude of the suboptimality of the $A_1$-optimal
weights relative to $\mathbf{w}^*(\kappa)$ is at most
\begin{equation}\label{eq:subopt}
    \frac{V(\mathbf{w}^{\mathrm{A_1}}) - V(\mathbf{w}^*(\kappa))}
         {V(\mathbf{w}^*(\kappa))} \;\leq\; \delta(\kappa),
\end{equation}
where $\delta(\kappa)$ is a function of the overlap
structure $\{\tau_{ij}\}$ and $\kappa$. In numerical experiments
with realistic crypto and equity timestamp structures, we find
$\delta(\kappa) < 1\%$ for all $\kappa \in [0,1]$ (see
Table~\ref{tab:subopt}), confirming that the $A_1$-optimal weights
are \emph{nearly efficient} in practice.
\end{remark}

\begin{table}[ht]
\centering
\caption{Variance reduction of $\mathbf{w}^*(\kappa)$ relative to
$A_1$-optimal weights, for a representative timestamp structure.
The $A_1$-optimal weights are full-variance optimal at $\rho=0$
and remain within 1\% of optimal for all $\rho \leq 1$.}
\label{tab:subopt}
\renewcommand{\arraystretch}{1.15}
\begin{tabular}{ccccc}
\toprule
$\rho$ & $V(\mathbf{w}^{\mathrm{A_1}})$ & $V(\mathbf{w}^*(\kappa))$
       & Reduction & Gap \\
\midrule
0.0  & 0.7443 & 0.7443 & 0.000\% & $A_1$-optimal = full optimal \\
0.3  & 0.7683 & 0.7683 & 0.003\% & negligible \\
0.5  & 0.8225 & 0.8222 & 0.038\% & negligible \\
0.7  & 0.9164 & 0.9134 & 0.327\% & small \\
0.9  & 1.0517 & 1.0443 & 0.697\% & small \\
1.0  & 1.1725 & 1.1610 & 0.975\% & $<1\%$ \\
\bottomrule
\end{tabular}
\end{table}

% ----------------------------------------------------------
\subsection{The Two-Step Estimator}\label{sec:twostep}
% ----------------------------------------------------------

The infeasibility of $\mathbf{w}^*(\kappa)$ motivates a
\emph{two-step} procedure in the spirit of feasible GLS,
using a first-step estimate of $\rho$ to construct an
approximation to the full-variance optimal weights.

\begin{definition}[Two-Step ZHY Estimator]\label{def:twostep}
The two-step ZHY estimator $\hat{\sigma}_{12}^{(2)}$ is defined
as follows.

\emph{Step~1.} Compute the $A_1$-optimal estimator
$\hat{\sigma}_{12}^{(1)}$ using weights
$w_{ij}^{(1)} = \tau_{ij}/(\tau_i^X\tau_j^Y)$.
Set
\begin{equation}\label{eq:kappa_hat}
    \hat\kappa
    \;=\;
    \left(\hat\rho^{(1)}\right)^2
    \;=\;
    \frac{\left(\hat\sigma_{12}^{(1)}\right)^2}
         {\hat\sigma_1^2\,\hat\sigma_2^2},
\end{equation}
where $\hat\sigma_1^2$ and $\hat\sigma_2^2$ are the realised
variances of assets $X$ and $Y$ over the same observation window.

\emph{Step~2.} Solve the KKT system~\eqref{eq:kkt} with
$\kappa = \hat\kappa$ to obtain weights
$\hat{\mathbf{w}}^* = \mathbf{w}^*(\hat\kappa)$.
Compute the final estimate
\begin{equation}\label{eq:twostep}
    \hat{\sigma}_{12}^{(2)}
    \;=\;
    \frac{\displaystyle\sum_{ij} R_i S_j\,\hat w_{ij}^*\,
          \mathbb{I}(\tau_{ij}>0)}
         {\displaystyle\sum_{ij} \tau_{ij}\,\hat w_{ij}^*}.
\end{equation}
\end{definition}

\begin{thm}[Two-Step Efficiency]\label{thm:twostep}
The two-step estimator $\hat{\sigma}_{12}^{(2)}$ is unbiased for
$\sigma_{12}$ and has exact conditional variance
\begin{equation}\label{eq:var_twostep}
    \VV\!\left(\hat{\sigma}_{12}^{(2)}\,\big|\,\tau,\,\hat\kappa\right)
    \;=\;
    \sigma_1^2\sigma_2^2\,A_1^{(2)}
    \;+\;
    \sigma_{12}^2\,A_2^{(2)},
\end{equation}
where $A_1^{(2)}$ and $A_2^{(2)}$ are the coefficients of
Theorem~\ref{thm:genHY} evaluated at the KKT weights
$\hat{\mathbf{w}}^*$.

Furthermore, conditionally on $\hat\kappa$, the two-step
estimator achieves the minimum variance within the class of
weighted ZHY estimators:
\begin{equation}\label{eq:dominance}
    \VV\!\left(\hat{\sigma}_{12}^{(2)}\,\big|\,\tau,\,\hat\kappa\right)
    \;\leq\;
    \VV\!\left(\hat{\sigma}_{12}(\mathbf{w})\,\big|\,\tau\right)
    \quad \text{for all weight matrices } \mathbf{w}.
\end{equation}
\end{thm}

\begin{proof}
Unbiasedness follows from~\eqref{eq:zhou_unbiased}, since the
two-step estimator is a weighted ZHY estimator and unbiasedness
holds for arbitrary non-negative weights satisfying
$T(\mathbf{w}) > 0$.

The exact variance formula~\eqref{eq:var_twostep} follows directly
from Theorem~\ref{thm:genHY} applied to the weights
$\hat{\mathbf{w}}^*$.

For the dominance statement~\eqref{eq:dominance}, fix
$\tau$ and $\hat\kappa$. The full conditional
variance~\eqref{eq:fullvar} is a strictly convex quadratic in
$\mathbf{w}$ under the normalisation $T(\mathbf{w}) = 1$. The
KKT system~\eqref{eq:kkt} with $\kappa = \hat\kappa$ characterises
the unique global minimum of $V(\mathbf{w})$ with
$\kappa = \hat\kappa$. Since $\hat{\mathbf{w}}^*$ solves
this system, it achieves the minimum, and~\eqref{eq:dominance}
follows.
\end{proof}

\begin{remark}[Computational cost]
The KKT system~\eqref{eq:kkt} is a $(K+1)\times(K+1)$ linear
system, where $K$ is the number of overlapping return-period pairs.
It is solved in $O(K^3)$ time using standard linear algebra
routines. Since the two-step estimator uses windowed processing
(Section~\ref{sec:HYrecursion}), $K$ is the number of overlapping
pairs within a single window, which is typically of order
$(\lambda_1 + \lambda_2)\,\Delta t$ where $\lambda_1, \lambda_2$
are the trade rates and $\Delta t$ is the window length.
For a one-hour window with rates of 50 and 5 trades per minute,
$K \approx 1{,}500$, giving a system of manageable size.
\end{remark}

\begin{remark}[Variance reduction]
The variance reduction of the two-step estimator relative to the
$A_1$-optimal estimator is exactly
\begin{equation}\label{eq:reduction}
    V\!\left(\mathbf{w}^{\mathrm{A_1}}\right)
    - V\!\left(\mathbf{w}^*(\hat\kappa)\right)
    \;=\;
    \sigma_1^2\sigma_2^2
    \bigl[A_1(\mathbf{w}^{\mathrm{A_1}}) - A_1(\mathbf{w}^*)\bigr]
    \;+\;
    \hat\kappa\,\sigma_1^2\sigma_2^2
    \bigl[A_2(\mathbf{w}^{\mathrm{A_1}}) - A_2(\mathbf{w}^*)\bigr].
\end{equation}
As Table~\ref{tab:subopt} shows, this quantity is small (typically
$< 1\%$) for realistic timestamp structures. The $A_1$-optimal
weights are therefore nearly efficient. The value of the two-step
estimator lies primarily in its theoretical completeness: it
establishes that the weighted ZHY class can achieve the
minimum variance for any given $\kappa$, and provides the
explicit constructive procedure for doing so.

The two-step estimator also serves as the natural candidate for
establishing an analogue of the Cram\'{e}r-Rao lower bound in the
asynchronous setting, as discussed in Section~\ref{sec:discussion}.
\end{remark}

\begin{remark}[Feasible two-step: consistency]
A sufficient condition for $\hat{\sigma}_{12}^{(2)}$ to have the
same asymptotic properties as the oracle two-step estimator
(using the true $\kappa$) is that $\hat\kappa \to \kappa$ in
probability as $\min(\lambda_1, \lambda_2) \to \infty$.
Since $\hat\kappa$ is a continuous function of the consistent
estimators $\hat\sigma_{12}^{(1)}$, $\hat\sigma_1^2$,
$\hat\sigma_2^2$, this follows from the consistency of the
$A_1$-optimal estimator established by
\citet{hayashi2005covariance}.
\end{remark}
